\documentclass[12pt]{article}

\input{../../../_assets/preambulo}

\begin{document}

    % 1. Foto de fondo
    % 2. Título
    % 3. Encabezado Izquierdo
    % 4. Color de fondo
    % 5. Coord x del titulo
    % 6. Coord y del titulo
    % 7. Fecha

    
	\input{../../../_assets/portada}
	\portadaExamenFotoDif{../../sapienza.jpg}{Computer Vision\\Examen II}{Computer Vision. Examen II}{MidnightBlue}{-8}{28}{2025-2026}{Joaquín Avilés de la Fuente}

    \begin{description}
        \item[Asignatura] Computer Vision.
        \item[Curso Académico] 2023-2024.
        \item[Grado] Grado en Informática.
        \item[Grupo] Grupo Único.
        \item[Profesor] Irene Amerini.
        \item[Descripción] Examen final escrito de Computer Vision.
        \item[Fecha] junio de 2024.
        \item[Duración] 60 min.
    
    \end{description}
    \newpage

    % ------------------------------------

    \begin{observacion}
        The total sum of all the exercises is 32 points, although the maximum mark of the exam is 30.
    \end{observacion}
    \begin{ejercicio}[8 points]
    Discuss the role of DFT (Discrete Fourier Transform) in image filtering and what information does it provide about an image. Provide an example of how DFT can be used to perform high-pass and low-pass filtering in the frequency domain.
    \end{ejercicio}

    \begin{ejercicio}[8 points]
        \begin{enumerate}
            \item[a.] Explain the Scale-Invariant Feature Transform (SIFT) algorithm. Outline the steps involved in detecting and describing features using SIFT How does it achieve scale and rotation invariance?
            \item[b.] Compare SIFT with the Speeded-Up Robust Features (SURF) algorithm. Discuss the differences in their approach and computational efficiency.
        \end{enumerate}
    \end{ejercicio}

    \begin{ejercicio}[8 points]
        \begin{enumerate}
            \item[a.] Explain the concept of fundamental matrix. How is it computed, and what information does it convey about the relationship between two images?
            \item[b.] What is the geometric meaning of the epipoles in the two images? How are they related to the fundamental matrix (algebraically)?
        \end{enumerate}
    \end{ejercicio}

    \begin{ejercicio}[8 points]
        Describe the process of computing a disparity map in stereo vision. How does the disparity map relate to depth information?
    \end{ejercicio}
\end{document}